% Options for packages loaded elsewhere
\PassOptionsToPackage{unicode}{hyperref}
\PassOptionsToPackage{hyphens}{url}
%
\documentclass[
  ignorenonframetext,
  aspectratio=169,
  russian,
]{beamer}
\usepackage{pgfpages}
\setbeamertemplate{caption}[numbered]
\setbeamertemplate{caption label separator}{: }
\setbeamercolor{caption name}{fg=normal text.fg}
\beamertemplatenavigationsymbolshorizontal
% Prevent slide breaks in the middle of a paragraph
\widowpenalties 1 10000
\raggedbottom
\setbeamertemplate{part page}{
  \centering
  \begin{beamercolorbox}[sep=16pt,center]{part title}
    \usebeamerfont{part title}\insertpart\par
  \end{beamercolorbox}
}
\setbeamertemplate{section page}{
  \centering
  \begin{beamercolorbox}[sep=12pt,center]{section title}
    \usebeamerfont{section title}\insertsection\par
  \end{beamercolorbox}
}
\setbeamertemplate{subsection page}{
  \centering
  \begin{beamercolorbox}[sep=8pt,center]{subsection title}
    \usebeamerfont{subsection title}\insertsubsection\par
  \end{beamercolorbox}
}
\AtBeginPart{
  \frame{\partpage}
}
\AtBeginSection{
  \ifbibliography
  \else
    \frame{\sectionpage}
  \fi
}
\AtBeginSubsection{
  \frame{\subsectionpage}
}

\usepackage{amsmath,amssymb}
\usepackage{iftex}
\ifPDFTeX
  \usepackage[T1]{fontenc}
  \usepackage[utf8]{inputenc}
  \usepackage{textcomp} % provide euro and other symbols
\else % if luatex or xetex
  \usepackage{unicode-math}
  \defaultfontfeatures{Scale=MatchLowercase}
  \defaultfontfeatures[\rmfamily]{Ligatures=TeX,Scale=1}
\fi
\usepackage{lmodern}
\usetheme[]{metropolis}
\usefonttheme{serif} % use mainfont rather than sansfont for slide text
\ifPDFTeX\else  
    % xetex/luatex font selection
    \setmainfont[Ligatures=TeX]{PT Serif}
    \setsansfont[Ligatures=TeX,Scale=MatchLowercase]{PT Sans}
    \setmonofont[Scale=MatchLowercase,Scale=0.9]{PT Mono}
\fi
% Use upquote if available, for straight quotes in verbatim environments
\IfFileExists{upquote.sty}{\usepackage{upquote}}{}
\IfFileExists{microtype.sty}{% use microtype if available
  \usepackage[]{microtype}
  \UseMicrotypeSet[protrusion]{basicmath} % disable protrusion for tt fonts
}{}
\usepackage{xcolor}
\newif\ifbibliography
\setlength{\emergencystretch}{3em} % prevent overfull lines
\setcounter{secnumdepth}{-\maxdimen} % remove section numbering


\providecommand{\tightlist}{%
  \setlength{\itemsep}{0pt}\setlength{\parskip}{0pt}}\usepackage{longtable,booktabs,array}
\usepackage{calc} % for calculating minipage widths
\usepackage{caption}
% Make caption package work with longtable
\makeatletter
\def\fnum@table{\tablename~\thetable}
\makeatother
\usepackage{graphicx}
\makeatletter
\newsavebox\pandoc@box
\newcommand*\pandocbounded[1]{% scales image to fit in text height/width
  \sbox\pandoc@box{#1}%
  \Gscale@div\@tempa{\textheight}{\dimexpr\ht\pandoc@box+\dp\pandoc@box\relax}%
  \Gscale@div\@tempb{\linewidth}{\wd\pandoc@box}%
  \ifdim\@tempb\p@<\@tempa\p@\let\@tempa\@tempb\fi% select the smaller of both
  \ifdim\@tempa\p@<\p@\scalebox{\@tempa}{\usebox\pandoc@box}%
  \else\usebox{\pandoc@box}%
  \fi%
}
% Set default figure placement to htbp
\def\fps@figure{htbp}
\makeatother

\IfFileExists{plex-otf.sty}{
  %% Full TeXlive
  % \usepackage[%
  %   % math,
  %   RM={Scale=0.94},SS={Scale=0.94},SScon={Scale=0.94},TT={Scale=MatchLowercase,FakeStretch=0.9},DefaultFeatures={Ligatures=Common}
  % ]{plex-otf}
}{
  %% TinyTeX
  \usepackage{libertine}
}

%%% Load theme
% https://deic.uab.cat/~iblanes/beamer_gallery/
\IfFileExists{beamerthemegotham.sty}{
  %% Full TeXlive
  \usetheme{gotham}
  \gothamset{
    numbering=totalpagenumber,
    parttocframe default=off,
    sectiontocframe default=off,
    subsectiontocframe default=off,
  }
}{
  %% TinyTeX
  \usetheme{Madrid}
}
\makeatletter
\@namedef{ver@babel.sty}{}
\makeatother
\usepackage{polyglossia}
\setdefaultlanguage{russian}
\setotherlanguage{english}
\metroset{progressbar=frametitle,sectionpage=progressbar,numbering=fraction}
\makeatletter
\beamer@ignorenonframefalse
\makeatother
\makeatletter
\@ifpackageloaded{caption}{}{\usepackage{caption}}
\AtBeginDocument{%
\ifdefined\contentsname
  \renewcommand*\contentsname{Содержание}
\else
  \newcommand\contentsname{Содержание}
\fi
\ifdefined\listfigurename
  \renewcommand*\listfigurename{Список иллюстраций}
\else
  \newcommand\listfigurename{Список иллюстраций}
\fi
\ifdefined\listtablename
  \renewcommand*\listtablename{Список таблиц}
\else
  \newcommand\listtablename{Список таблиц}
\fi
\ifdefined\figurename
  \renewcommand*\figurename{Рисунок}
\else
  \newcommand\figurename{Рисунок}
\fi
\ifdefined\tablename
  \renewcommand*\tablename{Таблица}
\else
  \newcommand\tablename{Таблица}
\fi
}
\@ifpackageloaded{float}{}{\usepackage{float}}
\floatstyle{ruled}
\@ifundefined{c@chapter}{\newfloat{codelisting}{h}{lop}}{\newfloat{codelisting}{h}{lop}[chapter]}
\floatname{codelisting}{Список}
\newcommand*\listoflistings{\listof{codelisting}{Листинги}}
\makeatother
\makeatletter
\makeatother
\makeatletter
\@ifpackageloaded{caption}{}{\usepackage{caption}}
\@ifpackageloaded{subcaption}{}{\usepackage{subcaption}}
\makeatother

\ifLuaTeX
\usepackage[bidi=basic]{babel}
\else
\usepackage[bidi=default]{babel}
\fi
\babelprovide[main,import]{russian}
\ifPDFTeX
\else
\babelfont{rm}[Ligatures=TeX]{PT Serif}
\fi
\babelprovide[import]{english}
% get rid of language-specific shorthands (see #6817):
\let\LanguageShortHands\languageshorthands
\def\languageshorthands#1{}
\usepackage{csquotes}
\usepackage{bookmark}

\IfFileExists{xurl.sty}{\usepackage{xurl}}{} % add URL line breaks if available
\urlstyle{same} % disable monospaced font for URLs
\hypersetup{
  pdftitle={Лабораторная работа №2},
  pdfauthor={Гайдук Софья Сергеевна},
  pdflang={ru-RU},
  hidelinks,
  pdfcreator={LaTeX via pandoc}}


\title{Лабораторная работа №2}
\subtitle{Презентация}
\author{Гайдук Софья Сергеевна}
\date{Invalid Date}
\institute{Российский университет дружбы народов, Москва, Россия}

\begin{document}
\frame{\titlepage}


\section{Информация}\label{ux438ux43dux444ux43eux440ux43cux430ux446ux438ux44f}

\begin{frame}{Докладчик}
\phantomsection\label{ux434ux43eux43aux43bux430ux434ux447ux438ux43a}
\begin{columns}[c]
\begin{column}{0.7\linewidth}
\begin{itemize}[<+->]
\tightlist
\item
  Гайдук Софья Сергеевна
\item
  Студент
\item
  Российский университет дружбы народов
\item
  \href{mailto:1032253645@pfur.ru}{\nolinkurl{1032253645@pfur.ru}}
\end{itemize}
\end{column}

\begin{column}{0.3\linewidth}
\end{column}
\end{columns}
\end{frame}

\begin{frame}{Цель}
\phantomsection\label{ux446ux435ux43bux44c}
Изучить идеологию и применение средств контроля версий. Освоить умения
по работе с git
\end{frame}

\begin{frame}{Задачи}
\phantomsection\label{ux437ux430ux434ux430ux447ux438}
Создать базовую конфигурацию для работы с git. Создать ключ SSH. Создать
ключ PGP. Настроить подписи git. Зарегистрироваться на Github. Создать
локальный каталог для выполнения заданий по предмету.
\end{frame}

\begin{frame}{Установка git и Установка gh}
\phantomsection\label{ux443ux441ux442ux430ux43dux43eux432ux43aux430-git-ux438-ux443ux441ux442ux430ux43dux43eux432ux43aux430-gh}
Для начала установим git и gh. В моём случае они уже установлены,
посмотрим версию

\begin{figure}[H]

{\centering \pandocbounded{\includegraphics[keepaspectratio]{image/image1.jpg}}

}

\caption{Установка git и gh}

\end{figure}%
\end{frame}

\begin{frame}{Базовая настройка git}
\phantomsection\label{ux431ux430ux437ux43eux432ux430ux44f-ux43dux430ux441ux442ux440ux43eux439ux43aux430-git}
Зададим имя для владельца репозитория.

\begin{figure}[H]

{\centering \pandocbounded{\includegraphics[keepaspectratio]{image/image2.jpg}}

}

\caption{Указание имени}

\end{figure}%
\end{frame}

\begin{frame}{Базовая настройка git}
\phantomsection\label{ux431ux430ux437ux43eux432ux430ux44f-ux43dux430ux441ux442ux440ux43eux439ux43aux430-git-1}
Зададим почту, на которую зарегистрирован аккаунт на github

\begin{figure}[H]

{\centering \pandocbounded{\includegraphics[keepaspectratio]{image/image2.jpg}}

}

\caption{Указание почты}

\end{figure}%
\end{frame}

\begin{frame}{Базовая настройка git}
\phantomsection\label{ux431ux430ux437ux43eux432ux430ux44f-ux43dux430ux441ux442ux440ux43eux439ux43aux430-git-2}
Настроим кодировку utf8 в выводе сообщений git

\begin{figure}[H]

{\centering \pandocbounded{\includegraphics[keepaspectratio]{image/image2.jpg}}

}

\caption{Настройка кодировки utf8}

\end{figure}%
\end{frame}

\begin{frame}{Базовая настройка git}
\phantomsection\label{ux431ux430ux437ux43eux432ux430ux44f-ux43dux430ux441ux442ux440ux43eux439ux43aux430-git-3}
Зададим имя начальной ветки, настроим параметры autocrlf и safecrlf

\begin{figure}[H]

{\centering \pandocbounded{\includegraphics[keepaspectratio]{image/image2.png}}

}

\caption{Настройка git}

\end{figure}%
\end{frame}

\begin{frame}{Создание ключа ssh}
\phantomsection\label{ux441ux43eux437ux434ux430ux43dux438ux435-ux43aux43bux44eux447ux430-ssh}
Создадим ключ RSA размером 4096 бит

\begin{figure}[H]

{\centering \includegraphics[width=\linewidth,height=0.6\textheight,keepaspectratio]{image/image3.png}

}

\caption{Создание ключа RSA}

\end{figure}%
\end{frame}

\begin{frame}{Создание ключа ssh}
\phantomsection\label{ux441ux43eux437ux434ux430ux43dux438ux435-ux43aux43bux44eux447ux430-ssh-1}
Теперь создадим ключ по алгоритму ed22519

\begin{figure}[H]

{\centering \includegraphics[width=\linewidth,height=0.6\textheight,keepaspectratio]{image/image4.png}

}

\caption{Создание ключа ed22519}

\end{figure}%
\end{frame}

\begin{frame}{Создание ключа pgp}
\phantomsection\label{ux441ux43eux437ux434ux430ux43dux438ux435-ux43aux43bux44eux447ux430-pgp}
Теперь создадим ключ gpg. Выбираем из предложенных вариантов первый тип
(RSA and RSA), размер ключа задаём 4096 бит и делаем срок действия ключа
неограниченным

\begin{figure}[H]

{\centering \includegraphics[width=\linewidth,height=0.6\textheight,keepaspectratio]{image/image5.png}

}

\caption{создание ключа pgp (1)}

\end{figure}%
\end{frame}

\begin{frame}{Создание ключа pgp}
\phantomsection\label{ux441ux43eux437ux434ux430ux43dux438ux435-ux43aux43bux44eux447ux430-pgp-1}
Вводим имя и адрес электронной почты.

\begin{figure}[H]

{\centering \includegraphics[width=\linewidth,height=0.6\textheight,keepaspectratio]{image/image6.png}

}

\caption{Создание ключа pgp (2)}

\end{figure}%
\end{frame}

\begin{frame}{Добавление PGP ключа в GitHub}
\phantomsection\label{ux434ux43eux431ux430ux432ux43bux435ux43dux438ux435-pgp-ux43aux43bux44eux447ux430-ux432-github}
Добавляем PGP ключа в GitHub

\begin{figure}[H]

{\centering \includegraphics[width=\linewidth,height=0.6\textheight,keepaspectratio]{image/image10.png}

}

\caption{добавление pgp ключей в Github}

\end{figure}%
\end{frame}

\begin{frame}{Настройка автоматических подписей коммитов git}
\phantomsection\label{ux43dux430ux441ux442ux440ux43eux439ux43aux430-ux430ux432ux442ux43eux43cux430ux442ux438ux447ux435ux441ux43aux438ux445-ux43fux43eux434ux43fux438ux441ux435ux439-ux43aux43eux43cux43cux438ux442ux43eux432-git}
Теперь производим настройку автоматических подписей

\begin{figure}[H]

{\centering \pandocbounded{\includegraphics[keepaspectratio]{image/image11.png}}

}

\caption{Настройка автоматических подписей коммитов git}

\end{figure}%
\end{frame}

\begin{frame}{Настройка автоматических подписей коммитов git}
\phantomsection\label{ux43dux430ux441ux442ux440ux43eux439ux43aux430-ux430ux432ux442ux43eux43cux430ux442ux438ux447ux435ux441ux43aux438ux445-ux43fux43eux434ux43fux438ux441ux435ux439-ux43aux43eux43cux43cux438ux442ux43eux432-git-1}
\begin{figure}[H]

{\centering \pandocbounded{\includegraphics[keepaspectratio]{image/image12.png}}

}

\caption{Настройка автоматических подписей коммитов git}

\end{figure}%
\end{frame}

\begin{frame}{Настройка автоматических подписей коммитов git}
\phantomsection\label{ux43dux430ux441ux442ux440ux43eux439ux43aux430-ux430ux432ux442ux43eux43cux430ux442ux438ux447ux435ux441ux43aux438ux445-ux43fux43eux434ux43fux438ux441ux435ux439-ux43aux43eux43cux43cux438ux442ux43eux432-git-2}
\begin{figure}[H]

{\centering \pandocbounded{\includegraphics[keepaspectratio]{image/image13.png}}

}

\caption{Настройка автоматических подписей коммитов git}

\end{figure}%
\end{frame}

\begin{frame}{Настройка gh}
\phantomsection\label{ux43dux430ux441ux442ux440ux43eux439ux43aux430-gh}
Нам нужно авторизоваться в github с помощью gh. Мы выбираем сайт для
авторизации (GitHub.com), после выбираем предпочитаемый протокол (SSH),
публичный SSH ключ (id\_rsa.pub), и имя для ключа (Sway). В качестве
способа авторизации выбираем браузер

\begin{figure}[H]

{\centering \includegraphics[width=\linewidth,height=0.45\textheight,keepaspectratio]{image/image14.png}

}

\caption{Авторизация в gh}

\end{figure}%
\end{frame}

\begin{frame}{Сознание репозитория курса на основе шаблона}
\phantomsection\label{ux441ux43eux437ux43dux430ux43dux438ux435-ux440ux435ux43fux43eux437ux438ux442ux43eux440ux438ux44f-ux43aux443ux440ux441ux430-ux43dux430-ux43eux441ux43dux43eux432ux435-ux448ux430ux431ux43bux43eux43dux430}
Теперь создаём рабочую директорию курса и переходим в неё

\begin{figure}[H]

{\centering \pandocbounded{\includegraphics[keepaspectratio]{image/image16.png}}

}

\caption{Создание рабочей директории и переход в неё}

\end{figure}%
\end{frame}

\begin{frame}{Сознание репозитория курса на основе шаблона}
\phantomsection\label{ux441ux43eux437ux43dux430ux43dux438ux435-ux440ux435ux43fux43eux437ux438ux442ux43eux440ux438ux44f-ux43aux443ux440ux441ux430-ux43dux430-ux43eux441ux43dux43eux432ux435-ux448ux430ux431ux43bux43eux43dux430-1}
Cоздаём репозиторий для лабораторных работ из шаблона

\begin{figure}[H]

{\centering \pandocbounded{\includegraphics[keepaspectratio]{image/17.png}}

}

\caption{Сознание репозитория курса}

\end{figure}%
\end{frame}

\begin{frame}{Сознание репозитория курса на основе шаблона}
\phantomsection\label{ux441ux43eux437ux43dux430ux43dux438ux435-ux440ux435ux43fux43eux437ux438ux442ux43eux440ux438ux44f-ux43aux443ux440ux441ux430-ux43dux430-ux43eux441ux43dux43eux432ux435-ux448ux430ux431ux43bux43eux43dux430-2}
И клонируем его к себе на компьютер

\begin{figure}[H]

{\centering \pandocbounded{\includegraphics[keepaspectratio]{image/18.png}}

}

\caption{Клонирование репозитория}

\end{figure}%
\end{frame}

\begin{frame}{Настройка каталога курса}
\phantomsection\label{ux43dux430ux441ux442ux440ux43eux439ux43aux430-ux43aux430ux442ux430ux43bux43eux433ux430-ux43aux443ux440ux441ux430}
Переходим в него с помощью cd и удаляем ненужные файлы (package.json) и
создаём необходимые каталоги, записав в файл COURSE строку os-intro (это
наш текущий курс) и прописываем make prepare для того, чтобы нужные нам
каталоги создались

\begin{figure}[H]

{\centering \includegraphics[width=\linewidth,height=0.45\textheight,keepaspectratio]{image/image19.png}

}

\caption{Удаление ненужных файлов и использование make}

\end{figure}%
\end{frame}

\begin{frame}{Настройка каталога курса}
\phantomsection\label{ux43dux430ux441ux442ux440ux43eux439ux43aux430-ux43aux430ux442ux430ux43bux43eux433ux430-ux43aux443ux440ux441ux430-1}
Теперь добавляем нашу папку для отправки

\begin{figure}[H]

{\centering \pandocbounded{\includegraphics[keepaspectratio]{image/image19.png}}

}

\caption{Использование git add}

\end{figure}%
\end{frame}

\begin{frame}{Настройка каталога курса}
\phantomsection\label{ux43dux430ux441ux442ux440ux43eux439ux43aux430-ux43aux430ux442ux430ux43bux43eux433ux430-ux43aux443ux440ux441ux430-2}
Делаем коммит, в котором указываем, что мы сделали структуру курса

\begin{figure}[H]

{\centering \includegraphics[width=\linewidth,height=0.6\textheight,keepaspectratio]{image/image20.png}

}

\caption{Использование git commit}

\end{figure}%
\end{frame}

\begin{frame}{Настройка каталога курса}
\phantomsection\label{ux43dux430ux441ux442ux440ux43eux439ux43aux430-ux43aux430ux442ux430ux43bux43eux433ux430-ux43aux443ux440ux441ux430-3}
И отправляем файлы на сервер GitHub с помощью команды push

\begin{figure}[H]

{\centering \pandocbounded{\includegraphics[keepaspectratio]{image/image23.png}}

}

\caption{Использование git push}

\end{figure}%
\end{frame}

\begin{frame}{Выводы}
\phantomsection\label{ux432ux44bux432ux43eux434ux44b}
Устаноли git, провели его первоначальную настройку, были созданы ключи
для авторизации и подписи, а также создан репозиторий курса из
предложенного шаблона
\end{frame}




\end{document}
